% Metodologia. Os números vêm de outputs/funnel_metrics*.json; [N] = ainda não calculado.
% O progresso fica em docs/checklist_artigo.md.
\section{Metodologia}

Esta pesquisa replica e adapta o protocolo de mineração de avaliações de aplicativos \parencite{gencnayebi2017} proposto por \textcite{alsebayel2025} para o domínio de finanças pessoais e educação financeira. O estudo investiga aplicativos comerciais disponíveis nas duas principais lojas móveis (Google Play Store e Apple App Store) em quatro mercados e três idiomas: Brasil (português), Estados Unidos (inglês), Espanha (espanhol) e México (espanhol). Para mitigar conflito de interesse, o aplicativo em desenvolvimento no contexto aplicado desta pesquisa não integra a amostra.

\subsection{Seleção dos aplicativos e coleta}

\textbf{Strings de busca.} Seguindo \textcite{alsebayel2025}, definimos previamente seis strings de busca (S1--S6), sem consultas por nomes de marca, preservando a mesma intenção de busca nas traduções para inglês e espanhol (Tabela~\ref{tab:strings}). Em 22 de novembro de 2026, cada string foi consultada individualmente na Google Play Store e na Apple App Store, registrando-se até os 30 primeiros resultados retornados por termo em cada loja e mercado.

\begin{table}[htbp]
\centering
\caption{Strings de busca aplicadas na Google Play Store e na Apple App Store por mercado e idioma.}
\label{tab:strings}
\small
\begin{tabular}{llll}
\toprule
\textbf{ID} & \textbf{Brasil (pt)} & \textbf{Estados Unidos (en)} & \textbf{Espanha e México (es)} \\
\midrule
S1 & \texttt{finanças} & \texttt{finance} & \texttt{finanzas} \\
S2 & \texttt{organizador de finanças} & \texttt{finance organizer} & \texttt{organizador de finanzas} \\
S3 & \texttt{educador de finanças} & \texttt{finance educator} & \texttt{educador de finanzas} \\
S4 & \texttt{aplicativos financeiros} & \texttt{financial apps} & \texttt{aplicaciones financieras} \\
S5 & \texttt{ia para finanças} & \texttt{ai for finance} & \texttt{ia para finanzas} \\
S6 & \texttt{organizador financeiro com ia} & \texttt{ai financial organizer} & \texttt{organizador financiero con ia} \\
\bottomrule
\end{tabular}
\end{table}

\textbf{Extração e deduplicação intraloja.} Na Google Play Store, a busca e a coleta de metadados utilizaram a biblioteca \texttt{google-play-scraper}\footnote{\url{https://github.com/JoMingyu/google-play-scraper}}, com recurso à página pública de busca quando a biblioteca falha ao ler resultados em destaque e complementação da data de atualização pela página pública do aplicativo quando o campo retorna ausente. Na Apple App Store, os metadados foram obtidos pela iTunes Search API oficial. Para cada ocorrência foram registrados identificador do aplicativo, título, desenvolvedor, categoria, número de avaliações, nota média, data da última atualização, descrição oficial, URL, faixa de instalações (apenas Google Play), termo de origem e posição no ranqueamento. Entradas duplicadas entre as seis strings dentro da mesma loja foram removidas pelo identificador único (\texttt{appId} na Google Play; \texttt{trackId} na App Store), preservando-se os termos de origem e a melhor posição obtida. Ao todo, as buscas retornaram 1.246 ocorrências (632 na Google Play e 614 na App Store).

\textbf{Critérios de elegibilidade.} Em conformidade com \textcite{alsebayel2025}, a triagem aplica critérios de elegibilidade em cascata: o aplicativo que falha em qualquer critério é excluído, registrando-se o primeiro motivo de exclusão de forma auditável:
\begin{enumerate}
    \item[(1)] \textit{Conflito de interesse:} exclui o aplicativo em desenvolvimento no contexto aplicado da pesquisa (0 exclusões).
    \item[(2)] \textit{Categoria:} o aplicativo deve pertencer à categoria ``Finanças'' (\textit{Finance} / \textit{Finanzas}); aplicativos da categoria ``Educação'' (\textit{Education} / \textit{Educación}) são admitidos somente quando o título ou o desenvolvedor indicam explicitamente educação financeira, excluindo-se cursos ou simulados voltados a certificações profissionais e \textit{trading}. Aplicativos de produtividade, jogos e demais categorias são excluídos.
    \item[(3)] \textit{Escopo funcional direto ao consumidor:} o aplicativo deve ser voltado a orçamento pessoal, controle de gastos ou educação financeira para indivíduos. Por meio de listas auditáveis aplicadas ao título, ao desenvolvedor e a sinais específicos da descrição oficial, excluem-se bancos e contas digitais, maquininhas de cartão e adquirentes, carteiras digitais e meios de pagamento ou remessa, plataformas de \textit{trading}, criptomoedas, corretoras e investimentos, ferramentas empresariais ou contábeis (B2B) e serviços de crédito, empréstimo ou adiantamento salarial.\footnote{Nos mercados internacionais, as listas de exclusão foram ampliadas com equivalentes em inglês e espanhol seguindo a mesma lógica aplicada ao Brasil.}
    \item[(4)] \textit{Presença de avaliações:} o aplicativo deve possuir ao menos uma avaliação de usuário registrada na loja.
    \item[(5)] \textit{Atualização recente:} o aplicativo deve apresentar data de atualização válida nos 12 meses anteriores à coleta (corte em 22 de novembro de 2025).
    \item[(6)] \textit{Disponibilidade nas duas lojas:} o aplicativo deve estar disponível simultaneamente na Google Play Store e na Apple App Store no mesmo mercado.
\end{enumerate}

Na triagem intraloja dos quatro mercados, foram excluídos 687 registros únicos (290 na Google Play e 397 na App Store), principalmente por categoria fora do escopo (224, sendo 180 em outras categorias e 44 em produtividade ou jogos), \textit{trading}, criptomoedas e investimentos (127), bancos, contas digitais e maquininhas (108), ausência de avaliações (93), carteiras digitais e pagamentos (55), ferramentas empresariais e crédito (48), desatualização superior a 12 meses (23) e sinais na descrição oficial (9).

\textbf{Verificação de disponibilidade e pareamento entre lojas.} Como o critério (6) exige presença do aplicativo nas duas lojas, e não necessariamente sua aparição simultânea entre os 30 primeiros resultados de busca em ambas, cada aplicativo elegível em apenas uma loja foi pesquisado pelo título na loja complementar, no mesmo mercado. Apenas correspondências com o mesmo nome-base de marca foram avaliadas pelas mesmas regras de elegibilidade em cascata, sendo incorporadas com origem rastreável (\texttt{availability\_check}) sem alterar as strings de busca. Em seguida, o pareamento 1:1 entre lojas operou em quatro camadas ordenadas por confiança: \texttt{exact} (título e desenvolvedor normalizados idênticos), \texttt{exact\_title} (título normalizado idêntico, admitindo variação de razão social), \texttt{normalized\_brand} (mesmo nome-base de marca e similaridade de desenvolvedor $\ge 0{,}75$) e \texttt{fuzzy\_review} (similaridade combinada $\ge 0{,}85$ ou mesma marca com desenvolvedor divergente), sendo esta última aceita apenas após confirmação manual. Nenhum par permaneceu pendente de revisão manual em qualquer mercado. Nos mercados internacionais, uma validação manual complementar descartou 4 aplicativos fora do escopo que haviam atravessado as regras lexicais (1 nos Estados Unidos, 1 na Espanha e 2 no México), resultando em 89 entradas aplicativo--mercado (67 aplicativos distintos, dos quais 12 figuram em mais de um mercado).

\textbf{Coleta das avaliações.} As avaliações públicas foram coletadas apenas para os aplicativos da amostra pareada em cada mercado. Na Google Play Store, a coleta empregou a biblioteca \texttt{google-play-scraper}, ordenada das avaliações mais recentes para as mais antigas. No Brasil, coletou-se todo o histórico acessível pela loja (223.688 avaliações de 40 aplicativos); nos mercados internacionais, coletaram-se diretamente as avaliações de 1, 2 e 3 estrelas na origem, limitadas a até 1.000 avaliações por nota e por aplicativo (até 3.000 por aplicativo). Na Apple App Store, diferentemente de \textcite{alsebayel2025} --- que a empregaram apenas na busca e triagem ---, a loja foi mantida como fonte de avaliações por decisão das autoras, visando ampliar a comparação entre plataformas. A coleta utilizou o \textit{feed} RSS público de avaliações da Apple (ordenado por \texttt{mostrecent}), cujo limite técnico é de 10 páginas de 50 entradas (no máximo as 500 avaliações mais recentes por aplicativo), concatenando título e corpo em um único texto. Ao todo, foram coletadas 278.046 avaliações brutas (262.678 na Google Play e 15.368 na App Store), detalhadas na Tabela~\ref{tab:coleta}. Em razão do limite de recência do \textit{feed} RSS e do recorte por estrelas na Google Play fora do Brasil, os volumes e horizontes temporais diferem entre plataformas e mercados. Além disso, as lojas da Espanha e do México podem retornar avaliações em espanhol oriundas de outros países, motivo pelo qual os dados são analisados por mercado de coleta e idioma, sem inferir o país de residência do autor.

\begin{table}[htbp]
\centering
\caption{Funil de busca, triagem, pareamento entre lojas e coleta bruta de avaliações ($N_0$) por mercado.}
\label{tab:coleta}
\footnotesize
\setlength{\tabcolsep}{3.5pt}
\begin{tabular}{lrrrrrrrrr}
\toprule
\textbf{Mercado} & \textbf{Ocorr.} & \textbf{Únicos} & \textbf{Eleg.} & \textbf{Eleg.} & \textbf{Pares} & \textbf{Amostra} & \textbf{Aval.} & \textbf{Aval.} & \textbf{Total} \\
 & \textbf{GP / AS} & \textbf{GP / AS} & \textbf{GP}$^a$ & \textbf{AS}$^a$ & \textbf{2 lojas} & \textbf{final}$^b$ & \textbf{GP} & \textbf{AS} & \textbf{($N_0$)} \\
\midrule
Brasil (pt)         & 180 / 170 &  98 / 147 & 58 & 64 & 40 & 40 & 223.688 &  5.467 & 229.155 \\
Estados Unidos (en) & 163 / 165 & 115 / 138 & 36 & 42 & 25 & 24 &  23.638 &  6.664 &  30.302 \\
Espanha (es)        & 145 / 138 & 101 / 123 & 24 & 35 & 12 & 11 &   7.441 &    901 &   8.342 \\
México (es)         & 144 / 141 & 106 / 122 & 34 & 29 & 16 & 14 &   7.911 &  2.336 &  10.247 \\
\midrule
\textbf{Total}      & \textbf{632 / 614} & \textbf{---} & \textbf{---} & \textbf{---} & \textbf{93} & \textbf{89}$^c$ & \textbf{262.678} & \textbf{15.368} & \textbf{278.046} \\
\bottomrule
\end{tabular}

\smallskip
\raggedright\scriptsize
$^a$Inclui os aplicativos elegíveis incorporados na verificação cruzada de disponibilidade (\texttt{availability\_check}): Brasil (+7 GP, +18 AS), Estados Unidos (+11 GP, +7 AS), Espanha (+0 GP, +7 AS) e México (+4 GP, +5 AS).\\
$^b$Após validação manual nos mercados internacionais (removidos 1 nos EUA, 1 na Espanha e 2 no México).\\
$^c$Corresponde a 67 aplicativos distintos; 12 aplicativos aparecem em mais de um mercado.
\end{table}

\textbf{Anonimização, mascaramento e redação de dados.} A anonimização ocorre desde a ingestão por política de lista de permissão (\textit{allowlist}): nomes de usuário, imagens de perfil, links e identificadores das avaliações são descartados na coleta, mantendo-se apenas a máscara do aplicativo, o texto, a nota, o idioma e a data. Cada aplicativo recebe uma máscara aleatória no formato \texttt{APP\_xxx}, sorteada de maneira não alfabética e mantida consistente entre lojas e mercados; a tabela de correspondência permanece em diretório privado fora do repositório e não é compartilhada, reduzindo o viés dos pesquisadores e evitando a atribuição de práticas a aplicativos específicos \parencite{alsebayel2025}. Na redação textual, regras determinísticas substituem e-mails (\texttt{[EMAIL]}), URLs (\texttt{[URL]}), CPFs (\texttt{[CPF]}), telefones (\texttt{[PHONE]}), sequências numéricas ambíguas (\texttt{[POSSIBLE\_PHONE\_OR\_NUMBER]}), identificadores com \texttt{@} (\texttt{[HANDLE]}) e nomes próprios autodeclarados em assinaturas ou saudações (\texttt{[NOME]}). Na preparação das cópias de trabalho e compartilhamento, uma expressão regular insensível a maiúsculas e acentos substitui as marcas dos aplicativos da amostra por \texttt{[APP]} --- desconsiderando termos coincidentes com palavras de uso comum para preservar a legibilidade --- e o modelo multilíngue de reconhecimento de entidades nomeadas \texttt{Davlan/bert-base-multilingual-cased-ner-hrl}\footnote{\url{https://huggingface.co/Davlan/bert-base-multilingual-cased-ner-hrl}} (limiar de confiança $\ge 0{,}6$, com expansão para palavra inteira e segmentação de textos longos em blocos de até 800 caracteres) substitui menções a pessoas (\texttt{[NOME]}), outras organizações (\texttt{[MARCA]}) e lugares (\texttt{[LOCAL]}). Ademais, as datas são truncadas ao dia (\texttt{AAAA-MM-DD}), removendo-se o horário exato, e uma verificação automatizada impede a gravação das tabelas públicas caso detecte marcas residuais de aplicativos, identificadores diretos ou datas com hora. Como a redação automática pode apresentar falhas e não garante anonimização absoluta, os trechos eventualmente citados no artigo são parafraseados para impedir a reidentificação dos avaliadores por busca reversa.

\textbf{Termos de serviço e considerações éticas.} Seguindo \textcite{krotov2020}, a coleta foi planejada considerando a legalidade, os termos das plataformas e o impacto sobre os sistemas e as pessoas envolvidas. Foram coletadas exclusivamente avaliações publicamente acessíveis, visíveis a qualquer visitante sem autenticação, e nenhum mecanismo de controle de acesso foi contornado. Na Google Play Store, as avaliações foram obtidas pela biblioteca \texttt{google-play-scraper}, adotando-se o mesmo escopo de dados públicos de \textcite{alsebayel2025}; na Apple App Store, utilizou-se o \textit{feed} RSS público de avaliações disponibilizado pela própria Apple. Para não sobrecarregar os serviços, as requisições foram espaçadas (1 segundo na Google Play e 3,2 segundos na App Store, com novas tentativas em intervalos crescentes em caso de falha) e o volume foi limitado (até 1.000 avaliações por nota e por aplicativo na Google Play nos mercados internacionais e até 500 avaliações por aplicativo na App Store). Os dados foram usados apenas para fins acadêmicos e não comerciais, e a coleta seguiu o princípio de minimização: nenhum dado do perfil do avaliador foi armazenado, os aplicativos foram mascarados e os resultados são relatados de forma agregada, com citações parafraseadas. Reconhecemos, contudo, que os Termos de Serviço das plataformas restringem o acesso automatizado e que sua aplicação à pesquisa acadêmica com dados públicos não é inequívoca; essa ambiguidade é declarada como limitação do estudo.
% TODO (orientação): confirmar esta redação antes da submissão.

\subsection{Pré-processamento}

Em conformidade com o protocolo adotado, o pré-processamento é aplicado \textbf{antes} da análise de sentimento e do filtro por estrelas, sempre sobre uma cópia dos dados, mantendo-se intacto o conjunto bruto coletado ($N_0$). Como unidade textual, definimos \textit{palavra} como sequência de caracteres delimitada por espaços em branco. O procedimento executa quatro operações sequenciais em cada mercado e loja: (1) remoção de avaliações com texto vazio após a supressão de espaços nas extremidades; (2) remoção de avaliações com três ou menos palavras ($\le 3$), por conterem informação insuficiente para a modelagem de tópicos; (3) geração de uma versão normalizada do texto em minúsculas com espaços padronizados, \textbf{sem} remoção de \textit{stopwords} nem \textit{stemming}, preservando o contexto semântico exigido pelos \textit{embeddings} na etapa de modelagem de tópicos; e (4) remoção de duplicatas exatas sobre o texto normalizado dentro de cada loja, retendo-se a ocorrência cronologicamente mais antiga. Para preservar sinais de ênfase lexical (como uso de letras maiúsculas) explorados pelo LeIA e pelo VADER, a análise de sentimento da etapa seguinte incide sobre o texto original das avaliações retidas em $N_1$.

No Brasil, das 223.688 avaliações coletadas na Google Play, foram removidas 158 avaliações vazias, 67.507 com três ou menos palavras e 5.739 duplicatas exatas, restando 150.284 avaliações ($N_1$); na App Store, das 5.467 coletadas, foram removidas 0 vazias, 315 curtas e 16 duplicatas, restando 5.136 avaliações (155.420 no total). Nos Estados Unidos, restaram 21.952 avaliações na Google Play (de 23.638 coletadas) e 6.451 na App Store (de 6.664), totalizando 28.403 avaliações. Na Espanha, restaram 6.876 avaliações na Google Play (de 7.441) e 866 na App Store (de 901), somando 7.742 avaliações; no México, restaram 7.327 na Google Play (de 7.911) e 2.195 na App Store (de 2.336), totalizando 9.522 avaliações. No conjunto dos quatro mercados, das 278.046 avaliações brutas ($N_0$), foram removidas 170 avaliações vazias, 70.999 com três ou menos palavras e 5.790 duplicatas exatas (76.959 remoções), resultando em 201.087 avaliações pré-processadas ($N_1$), conforme detalhado na Tabela~\ref{tab:preprocessamento}.

\begin{table}[htbp]
\centering
\caption{Funil de pré-processamento ($N_0 \to N_1$) por mercado e loja.}
\label{tab:preprocessamento}
\small
\begin{tabular}{lrrrrr}
\toprule
\textbf{Mercado / Loja} & \textbf{Brutas ($N_0$)} & \textbf{Vazias} & \textbf{$\le 3$ pal.} & \textbf{Duplicatas} & \textbf{Restantes ($N_1$)} \\
\midrule
Brasil --- Google Play         & 223.688 & 158 & 67.507 & 5.739 & 150.284 \\
Brasil --- App Store           &   5.467 &   0 &    315 &    16 &   5.136 \\
Estados Unidos --- Google Play &  23.638 &   2 &  1.664 &    20 &  21.952 \\
Estados Unidos --- App Store   &   6.664 &   0 &    212 &     1 &   6.451 \\
Espanha --- Google Play        &   7.441 &   5 &    556 &     4 &   6.876 \\
Espanha --- App Store          &     901 &   0 &     33 &     2 &     866 \\
México --- Google Play         &   7.911 &   5 &    575 &     4 &   7.327 \\
México --- App Store           &   2.336 &   0 &    137 &     4 &   2.195 \\
\midrule
\textbf{Total}                 & \textbf{278.046} & \textbf{170} & \textbf{70.999} & \textbf{5.790} & \textbf{201.087} \\
\bottomrule
\end{tabular}
\end{table}

\subsection{Análise de sentimento e filtro de criticidade}

Para identificar relatos críticos sobre o conjunto pré-processado ($N_1$), aplicamos classificadores adequados ao idioma de cada mercado: o LeIA\footnote{\url{https://github.com/rafjaa/LeIA}}, adaptação léxica do VADER para o português \parencite{pereira2021}, no Brasil; o VADER \parencite{hutto2014}, nos Estados Unidos; e o \textit{pysentimiento} \parencite{perez2021}, baseado no modelo \texttt{robertuito-sentiment-analysis}, na Espanha e no México. No LeIA e no VADER, adotaram-se os limiares padrão do índice \textit{compound}: valores $\ge 0{,}05$ indicam sentimento positivo, valores $\le -0{,}05$ indicam sentimento negativo e valores no intervalo intermediário ($-0{,}05 < \textit{compound} < 0{,}05$) indicam sentimento neutro. No \textit{pysentimiento}, a decisão utiliza diretamente o rótulo previsto (\texttt{POS}, \texttt{NEU} ou \texttt{NEG}).

Seguindo \textcite{alsebayel2025}, aplicamos uma regra de \textbf{interseção}: foram retidas somente as avaliações com nota igual ou inferior a três estrelas ($\le 3$) \textbf{e} sentimento negativo ou neutro (\textit{compound} $< 0{,}05$ em português e inglês; rótulos \texttt{NEG} ou \texttt{NEU} em espanhol), descartando-se as avaliações classificadas como positivas. A Tabela~\ref{tab:sentimento} apresenta a distribuição de polaridade sobre $N_1$, o volume que seria retido por um filtro baseado apenas na nota em estrelas e o total de \textit{avaliações críticas} ($N_2$) obtido pela combinação. No conjunto dos quatro mercados, o filtro apenas por estrelas reteria 73.675 avaliações (excluindo 127.412 com nota $> 3$); a interseção com o sentimento descartou ainda 21.912 avaliações positivas, resultando em 51.763 \textit{avaliações críticas} ($N_2$): 24.336 no Brasil (22.316 Google Play e 2.020 App Store), 13.586 nos Estados Unidos (12.188 e 1.398), 6.587 na Espanha (6.305 e 282) e 7.254 no México (6.735 e 519).\footnote{Ressalta-se que, como a coleta na Google Play dos mercados internacionais já filtrou notas de 1 a 3 estrelas na origem, a proporção de avaliações negativas nesses recortes não é diretamente comparável à do Brasil ou da App Store; além disso, nos aplicativos presentes tanto na Espanha quanto no México, há sobreposição de avaliações em espanhol na Google Play.}

\begin{table}[htbp]
\centering
\caption{Distribuição de sentimento em $N_1$ e filtro de criticidade ($N_1 \to N_2$) por mercado e loja.}
\label{tab:sentimento}
\footnotesize
\setlength{\tabcolsep}{3.5pt}
\begin{tabular}{lrrrrrrr}
\toprule
\textbf{Mercado / Loja} & \textbf{$N_1$} & \textbf{Pos. (\%)} & \textbf{Neu. (\%)} & \textbf{Neg. (\%)} & \textbf{Nota $\le 3$} & \textbf{Pos. rem.} & \textbf{Críticas ($N_2$)} \\
\midrule
Brasil --- GP         & 150.284 & 68{,}50 & 10{,}52 & 20{,}98 & 31.507 &  9.191 & 22.316 \\
Brasil --- AS         &   5.136 & 50{,}31 &  7{,}50 & 42{,}19 &  2.595 &    575 &  2.020 \\
EUA --- GP$^a$        &  21.952 & 44{,}48 & 18{,}43 & 37{,}09 & 21.952 &  9.764 & 12.188 \\
EUA --- AS            &   6.451 & 71{,}48 &  7{,}81 & 20{,}71 &  2.574 &  1.176 &  1.398 \\
Espanha --- GP$^a$    &   6.876 &  8{,}30 & 25{,}12 & 66{,}58 &  6.876 &    571 &  6.305 \\
Espanha --- AS        &     866 & 55{,}77 & 11{,}43 & 32{,}79 &    299 &     17 &    282 \\
México --- GP$^a$     &   7.327 &  8{,}08 & 25{,}07 & 66{,}85 &  7.327 &    592 &  6.735 \\
México --- AS         &   2.195 & 62{,}60 & 11{,}75 & 25{,}65 &    545 &     26 &    519 \\
\midrule
\textbf{Total}        & \textbf{201.087} & \textbf{---} & \textbf{---} & \textbf{---} & \textbf{73.675} & \textbf{21.912} & \textbf{51.763} \\
\bottomrule
\end{tabular}

\smallskip
\raggedright\scriptsize
$^a$Coletadas na Google Play apenas para notas de 1 a 3 estrelas; portanto, todas as avaliações em $N_1$ possuem nota $\le 3$.
\end{table}

Para auditar a qualidade da classificação automática, foram geradas amostras aleatórias reprodutíveis (\texttt{seed = 42}) contendo até 100 avaliações retidas como críticas e até 100 avaliações com nota $\le 3$ descartadas por sentimento positivo em cada mercado e loja. A taxa de erro verificada nessa checagem manual foi de [N]\%.
% TODO: preencher [N]% com a taxa de erro após a revisão humana dos arquivos checagem_manual_sentimento.csv.
Reconhecemos que relatos relevantes podem ter sido excluídos por discrepância entre o sentimento atribuído automaticamente e a nota em estrelas, sobretudo em textos irônicos ou mistos.

\subsection{Modelagem de tópicos}
% TODO (Etapa 3).

\subsection{Inspeção manual dos tópicos}
% TODO (Etapa 4) + tabela do funil.

\subsection{Análise temática reflexiva}
% TODO (Etapa 5).
